\section{Limitations}\label{sec:limitations}

Motivo occupies a deliberate niche: offline, rule-based expansion into Standard MIDI Files.
That focus enables determinism and a compact implementation, but it also defines what the system does \emph{not} do.
The limitations below are stated honestly so that the contributions in Section~\ref{sec:contributions} are not read as
claims about a general-purpose digital audio workstation or live performance environment.

\subsection{Expressive MIDI and sound design}\label{subsec:lim-expressive-midi}

The intermediate representation stores pitch, start beat, duration, and a fixed default velocity for each note event.
The source language does not expose per-note velocity literals, continuous controllers, pitch bend, or filter-related
MIDI data.
Composers cannot shape dynamics, timbral sweeps, or synthesizer filter parameters from Motivo source code; those
decisions remain downstream in a DAW or sampler.
Motivo emits straight note-on and note-off messages on General MIDI program assignments rather than driving a synthesis
graph, an audio engine, or effect chains.
The Studio playback layer uses SoundFont rendering for auditioning, but that path is separate from the compiler output
and does not extend the language's expressive range.

\subsection{Rigid instruments and single-file scope}\label{subsec:lim-instruments-modularity}

Instrument identity is fixed at compile time through a small closed set of keywords (\texttt{piano}, \texttt{guitar},
\texttt{bass}, \texttt{violin}, \texttt{drums}) mapped to General MIDI program numbers or the percussion channel.
There is no user-defined instrument registry, custom sound-font binding, or per-project timbre configuration in the
language.
All declarations for a composition must live in one source file: Motivo has no \texttt{import} or module mechanism,
no cross-file pattern libraries, and no project-level packaging beyond what Motivo Studio stores as individual
documents.
Reusable ideas are limited to patterns defined within the same translation unit, which constrains large-scale
modularity compared with multi-file software projects or DAW session hierarchies.

\subsection{Workflow fit versus DAW note entry}\label{subsec:lim-daw-workflow}

For exploratory scratch composition---placing a few notes, nudging timing by ear, or sketching a harmonic idea on a
piano roll---a conventional DAW remains faster and more direct.
Motivo requires the composer to think in terms of rules, loops, and typed declarations before hearing a result, and
each edit passes through compile diagnostics rather than immediate grid manipulation.
The language is therefore better suited to algorithmically structured repetition and variation than to freehand
micro-editing of individual events.
This is an event-placement and compile-time model, not an active music-production workflow with real-time manipulation,
automation lanes, mixing, or arrangement editing on a timeline.

\subsection{Language coverage and evaluation scope}\label{subsec:lim-language-eval}

Several features common in production-oriented tools are absent or only partially modeled.
Key-signature declarations and harmonic context are not part of the current grammar; accidentals are written as
literal pitches.
The pattern system is static: parameters accept concrete evaluated types, but there is no higher-order pattern
algebra or shipped standard library of musical transformations.
The lowerer enforces a hard cap of 20,000 note events to prevent runaway programs, which bounds the largest scores the
implementation can emit today.

The evaluation in Chapter~\ref{ch:evaluation} is technical rather than user-centered.
It reports compilation time, source-to-event expansion, MIDI size, and automated test counts on synthetic fixtures and
a single development machine.
It does not measure whether composers prefer Motivo over a DAW, whether beginners learn the syntax quickly, or whether
pieces written in Motivo are musically stronger.
Those questions would require a different study design with participants and qualitative instruments; this thesis does
not claim them.
